\section{Fundamentos teóricos}
\label{sec:fundamentos}

% TODO: desarrollar cada subsección con citas
\subsection{Aprendizaje automático}

El aprendizaje automático es una rama de la inteligencia artificial que permite a los modelos aprender patrones a partir de datos sin ser programados explícitamente \parencite{hastie2009elements}. Entre los algoritmos más utilizados para clasificación se encuentran los árboles de decisión, los bosques aleatorios y las máquinas de vectores de soporte \parencite{geron2019hands}.

\subsection{Análisis espacial del delito}

La criminología ambiental estudia la relación entre el delito y el espacio, señalando que los eventos delictivos no se distribuyen de manera uniforme sino que se concentran en puntos calientes o ``hot spots'' \parencite{eck2005mapping}. Estos patrones pueden identificarse mediante técnicas de estadística espacial y sistemas de información geográfica (SIG).

\subsection{Patrones espaciales de criminalidad}

Los patrones espaciales de criminalidad describen la forma en que los delitos se agrupan, se repiten o se desplazan en el territorio. Su identificación requiere transformar los datos de incidentes en variables espaciales como densidad, cercanía o vecindario, y luego aplicar métodos de clasificación \parencite{ratcliffe2008intelligence}.